\scanpage{122}
2) \(F\) \deleted{\(\unsure{=Q_{lX}}\)} strictement spécial
\inserted{de poids \(j\)}, \(X\) régulier.

\diagram{\[
\begin{tikzcd}[name=diag122,execute at end picture={
\draw[rounded corners] (-.35,.95) rectangle (1.5,-.15);
}]
X'\arrow[r,hook,no head,"f'"]\arrow[d,"p"']&
Y'\arrow[d,"q"]\\
X\arrow[r,hook,"f"']&\deleted{X}\,\unsure{Y}
\end{tikzcd}
\]}
Cas très spécial
\[
F=R^{\unsure{d|j}}p_*(Q_{lX'}),
\]
\margin{\(p\) de dim relative \(j\)}

On a alors
\[
R^\alpha(fp)_*(Q_{lX'})
\Longleftarrow E_2^{pq}=R^pf_*(R^q\unsure{g_*|p_*}(Q_l))
\]
suite spectrale qui \uline{dégénère} par Deligne,
d'où
\[
\unsure{R^if_*}(F)=R^if_*(R^j\unsure{g_*|p_*}(Q_{lX'}))
\]
est un \uline{morceau} de
\[
R^{i+j}(fp)_*(Q_{lX'})
\]
Or
\[
R^\alpha(fp)_*(Q_{lX'})
\Longleftarrow \unsure{E}_2^{pq}=R^pq_*(R^qf'_*(Q_{lX'}))
\]
Or \(Y'\) est de dim \(\deleted{\unsure{n+j}}\,n'\le\boxed{n+j}\)

donc \(R^qf'_*(Q_{lX'})\) est effectif de poids
\[
\le2\operatorname{Inf}(\unsure{2q|q},\unsure{2n+2j|n+j})
\]
de niveau
\[
\le\operatorname{Inf}(q,n+j-1)
\]
\(q\) est de dim relative \(\le\)

\deleted{On prend \(p+q=i+j\)}

\deleted{\unsure{a priori}} \(E_2^{pq}\) de poids \(\le\)
% archive foliation: 121
